%% =================================================================
%% Kingma, Ba. "Adam: A Method for Stochastic Optimization."
%% ICLR 2015. arXiv:1412.6980v9.
%%
%% Reconstruction of page 5 (printed p. 5) as study notes.
%% Generated by: reproduce-all-pages-basic-tex (batch 1-15)
%%
%% Structure: Sec 5 Related Work continues (SFO/NGD blurb + RMSProp
%% paragraph + AdaGrad paragraph with inline derivation showing Adam
%% with β1=0 reduces to AdaGrad) + Sec 6 EXPERIMENTS + Sec 6.1
%% EXPERIMENT: LOGISTIC REGRESSION (cut off at end).
%% No new numbered equations.
%% =================================================================

\documentclass[11pt]{article}

\usepackage[margin=1in]{geometry}
\usepackage{amsmath, amssymb}
\usepackage{mathrsfs}
\usepackage{bm}
\usepackage{xcolor}
\usepackage{fancyhdr}

\pagestyle{fancy}
\fancyhf{}
\lhead{\small Published as a conference paper at ICLR 2015}
\rhead{}
\cfoot{\small 5 $\to$ \arabic{page}}
\renewcommand{\headrulewidth}{0.4pt}

\setcounter{page}{1}
\setcounter{equation}{5}

\begin{document}

\noindent
from first-order information. The Sum-of-Functions Optimizer
(SFO) (Sohl-Dickstein et al., 2014) is a quasi-Newton method
based on minibatches, but (unlike Adam) has memory requirements
linear in the number of minibatch partitions of a dataset, which
is often infeasible on memory-constrained systems such as a GPU.
Like natural gradient descent (NGD) (Amari, 1998), Adam employs
a preconditioner that adapts to the geometry of the data, since
$\widehat{v}_t$ is an approximation to the diagonal of the Fisher
information matrix (Pascanu \& Bengio, 2013); however, Adam's
preconditioner (like AdaGrad's) is more conservative in its
adaption than vanilla NGD by preconditioning with the square root
of the inverse of the diagonal Fisher information matrix
approximation.

\vspace{0.5em}

\noindent\textbf{RMSProp:} An optimization method closely related
to Adam is RMSProp (Tieleman \& Hinton, 2012). A version with
momentum has sometimes been used (Graves, 2013). There are a few
important differences between RMSProp with momentum and Adam:
RMSProp with momentum generates its parameter updates using a
momentum on the rescaled gradient, whereas Adam updates are
directly estimated using a running average of first and second
moment of the gradient. RMSProp also lacks a bias-correction
term; this matters most in case of a value of $\beta_2$ close to
$1$ (required in case of sparse gradients), since in that case
not correcting the bias leads to very large stepsizes and often
divergence, as we also empirically demonstrate in section~6.4.

\vspace{0.5em}

\noindent\textbf{AdaGrad:} An algorithm that works well for
sparse gradients is AdaGrad (Duchi et al., 2011). Its basic
version updates parameters as
$\theta_{t+1} = \theta_t - \alpha \cdot g_t / \sqrt{\sum_{i=1}^{t} g_i^2}$.
Note that if we choose $\beta_2$ to be infinitesimally close to
$1$ from below, then
$\lim_{\beta_2 \to 1} \widehat{v}_t = t^{-1} \cdot \sum_{i=1}^{t} g_i^2$.
AdaGrad corresponds to a version of Adam with $\beta_1 = 0$,
infinitesimal $(1 - \beta_2)$ and a replacement of $\alpha$ by an
annealed version $\alpha_t = \alpha \cdot t^{-1/2}$, namely
$\theta_t - \alpha \cdot t^{-1/2} \cdot \widehat{m}_t / \sqrt{\lim_{\beta_2 \to 1} \widehat{v}_t}
 = \theta_t - \alpha \cdot t^{-1/2} \cdot g_t / \sqrt{t^{-1} \cdot \sum_{i=1}^{t} g_i^2}
 = \theta_t - \alpha \cdot g_t / \sqrt{\sum_{i=1}^{t} g_i^2}$.
Note that this direct correspondence between Adam and AdaGrad
does not hold when removing the bias-correction terms; without
bias correction, like in RMSProp, a $\beta_2$ infinitesimally
close to $1$ would lead to infinitely large bias, and infinitely
large parameter updates.

\vspace{0.5em}

\section*{6 \textsc{Experiments}}

To empirically evaluate the proposed method, we investigated
different popular machine learning models, including logistic
regression, multilayer fully connected neural networks and deep
convolutional neural networks. Using large models and datasets,
we demonstrate Adam can efficiently solve practical deep learning
problems.

We use the same parameter initialization when comparing different
optimization algorithms. The hyper-parameters, such as learning
rate and momentum, are searched over a dense grid and the results
are reported using the best hyper-parameter setting.

\subsection*{6.1 \textsc{Experiment: Logistic regression}}

We evaluate our proposed method on L2-regularized multi-class
logistic regression using the MNIST dataset. Logistic regression
has a well-studied convex objective, making it suitable for
comparison of different optimizers without worrying about local
minimum issues. The stepsize $\alpha$ in our logistic regression
experiments is adjusted by $1/\sqrt{t}$ decay, namely
$\alpha_t = \frac{\alpha}{\sqrt{t}}$ that matches with our
theoratical prediction from section~4. The logistic regression
classifies the class label directly on the 784 dimension image
vectors. We compare Adam to accelerated SGD with Nesterov
momentum and Adagrad using minibatch size of 128. According to
Figure~1, we found that the Adam yields similar convergence as
SGD with momentum and both converge faster than Adagrad.

As discussed in (Duchi et al., 2011), Adagrad can efficiently
deal with sparse features and gradients as one of its main
theoretical results whereas SGD is low at learning rare features.
Adam with $1/\sqrt{t}$ decay on its stepsize should theoretically
match the performance of Adagrad. We examine the sparse feature
problem using IMDB movie review dataset from (Maas et al., 2011).
We pre-process the IMDB movie reviews into bag-of-words (BoW)
feature vectors including the first 10,000 most frequent words.
The 10,000 dimension BoW feature vector for each review is highly
sparse. As suggested in (Wang \& Manning, 2013), 50\% dropout
noise can be applied to the BoW features during

% [page ends here — continues on page 6]

\end{document}
